% Größen und Einheiten – bank/einheiten/ (zusammenbau v0.4, Heft)
\einheitenkopf[t7]{Größen und Einheiten}
\setcounter{aufgabe}{52}

% Anlauf (ohne Prüfkennung)
\begin{aufgabe}{Länge, Masse und Geld umrechnen – Anlauf}
\begin{teile}
\teil Welche Einheit ist größer? Kreuze an.\\ \kreuz{Meter}\\ \kreuz{Millimeter}
\teil $7$ m – wie viel dm? \leerfeld[dm]
\teil $30$ mm – wie viel cm? \leerfeld[cm]
\end{teile}
\end{aufgabe}

\begin{aufgabe}{Länge, Masse und Geld umrechnen – Prüfungsaufgaben}
\begin{teile}
\teil Setze $<$, $>$ oder $=$ ein: $0{,}04$ m \leerfeld $40$ cm \hfill (P10 2019 OS)
\teil Setze $<$, $=$ oder $>$ ein: $2{,}7$ m \leerfeld $27$ cm \hfill (P10 2026 FOR)
\teil Setze $<$, $>$ oder $=$ ein: $0{,}07$ kg \leerfeld $70$ g \hfill (P10 2019 OS)
\teil Ein Fernseher hat eine Bildschirmdiagonale von $55$ Zoll; $1$ Zoll $\approx 2{,}54$ cm. Weise nach, dass die Diagonale etwa $1{,}40$ m lang ist. \hfill (P10 2016 OS)
\teil Ein Volkslauf ist $6$ Meilen lang; $1$ Meile $\approx 1{,}61$ km. Weise nach, dass die Strecke knapp $9{,}7$ km lang ist. \hfill (P10 2016 OS)
\end{teile}
\end{aufgabe}

% Anlauf (ohne Prüfkennung)
\begin{aufgabe}{Zeit umrechnen und Zeitspannen berechnen – Anlauf}
\begin{teile}
\teil Der Film beginnt um 19:15 Uhr. Zeitpunkt oder Zeitspanne? Kreuze an.\\ \kreuz{Zeitpunkt}\\ \kreuz{Zeitspanne}
\teil $3$ min – wie viele s? \leerfeld[s]
\teil $240$ s – wie viele min? \leerfeld[min]
\end{teile}
\end{aufgabe}

\begin{aufgabe}{Zeit umrechnen und Zeitspannen berechnen – Prüfungsaufgaben}
\begin{teile}
\teil Eine Radtour dauert $5{,}5$ h. Wie viele Minuten sind das? \hfill (P10 2018 OS) \\ \leerfeld[min]
\teil Eine Schicht dauert $6{,}5$ Stunden. Gib die Dauer in Minuten an. \hfill (P10 2025 OS) \\ \leerfeld[min]
\teil Wie viele Minuten sind $2{,}25$ Stunden? Kreuze an. \hfill (P10 2024 OS)\\ \kreuz{$125$ min}\\ \kreuz{$135$ min}\\ \kreuz{$145$ min}\\ \kreuz{$155$ min}
\teil Ein Zug fährt um 9:47 Uhr ab und ist $3$ h $25$ min unterwegs. Wann kommt er an? \hfill (P10 2021 OS) \\ \leerfeld Uhr
\teil Ein Fernbus fährt um 16:52 Uhr ab, die Fahrt dauert $1$ h $43$ min. Wann kommt er an? \hfill (P10 2021 OS) \\ \leerfeld Uhr
\end{teile}
\end{aufgabe}

% Anlauf (ohne Prüfkennung)
\begin{aufgabe}{Flächen- und Volumeneinheiten umrechnen – Anlauf}
\begin{teile}
\teil Quadratmeter und Quadratdezimeter – welche Umrechnungszahl? Kreuze an.\\ \kreuz{zehn}\\ \kreuz{hundert}\\ \kreuz{tausend}
\teil $3$ m² – wie viel dm²? \leerfeld dm²
\teil $4$ m² – wie viel cm²? \leerfeld[cm²]
\end{teile}
\end{aufgabe}

\begin{aufgabe}{Flächen- und Volumeneinheiten umrechnen – Prüfungsaufgaben}
\begin{teile}
\teil Eine Flasche enthält $1{,}25$ l Saft. Wie viele Gläser zu $200$ ml kann man damit ganz füllen? \hfill (P10 2022 OS) \\ \leerfeld Gläser
\teil Ein Kanister enthält $4{,}2$ l Wasser. Wie viele Trinkflaschen zu $750$ ml kann man damit ganz füllen? \hfill (P10 2022 OS) \\ \leerfeld Flaschen
\teil In einem Topf sind $2{,}7$ l Suppe. Wie viele volle Kellen zu $250$ ml kann man ausgeben? \hfill (P10 2022 OS) \\ \leerfeld Kellen
\teil Eine Pumpe fördert $12\,500$ l pro Stunde. Wie lange braucht sie, um einen Teich mit $75$ m³ Wasser zu füllen? \hfill (P10 2017 OS) \\ \leerfeld[h]
\teil Ein Löschfahrzeug pumpt $800$ l pro Minute in einen Löschwasserbehälter mit $36$ m³. Wie lange dauert das Füllen? \hfill (P10 2017 OS) \\ \leerfeld[min]
\end{teile}
\end{aufgabe}

% Anlauf (ohne Prüfkennung)
\begin{aufgabe}{Mit Größen rechnen im Sachzusammenhang – Anlauf}
\begin{teile}
\teil Von einer $1{,}2$ m langen Leiste werden $45$ cm abgesägt; gefragt ist der Rest in cm. Welche Angabe musst du vorher umrechnen? Kreuze an.\\ \kreuz{$1{,}2$ m}\\ \kreuz{$45$ cm}
\teil $1{,}4$ km $+ 650$ m – wie viel km? \leerfeld[km]
\teil Ein Modellhaus ist $0{,}08$ m hoch. Gib die Höhe in einer handlicheren Einheit an. \leerfeld \leerfeld
\end{teile}
\end{aufgabe}

\begin{aufgabe}{Mit Größen rechnen im Sachzusammenhang – Prüfungsaufgaben}
\begin{teile}
\teil Eine Messingkugel wiegt $3\,500$ g; Messing hat die Dichte $8{,}4$ g/cm³. Berechne das Volumen der Kugel. Runde auf eine Stelle nach dem Komma. \hfill (P10 2016 OS) \\ \leerfeld[cm³]
\teil Ein Anglergewicht aus Blei wiegt $1\,800$ g; Blei hat die Dichte $11{,}3$ g/cm³. Berechne sein Volumen. Runde auf eine Stelle nach dem Komma. \hfill (P10 2016 OS) \\ \leerfeld[cm³]
\teil Ein Zinnbecher wiegt $620$ g; Zinn hat die Dichte $7{,}3$ g/cm³. Berechne das Volumen des Zinns. Runde auf eine Stelle nach dem Komma. \hfill (P10 2016 OS) \\ \leerfeld[cm³]
\teil Ein Vogelhaus hat zwei Dachflächen zu je $0{,}18$ m² und vier Wände zu je $0{,}12$ m². Alles wird zweimal gestrichen; $1$ l Farbe reicht für $8$ m². Es gibt Dosen mit $125$ ml, $250$ ml und $750$ ml. Lea meint, die kleinste Dose reicht, Max will die mittlere kaufen. Wer hat recht? Begründe. \hfill (P10 2015 OS)
\teil Ein Gartenschrank hat vier Außenwände zu je $0{,}9$ m², die zweimal lasiert werden; $1$ l Lasur reicht für $12$ m². Es gibt Dosen mit $500$ ml, $750$ ml und $1$ l. Tom will die 1-l-Dose, Ina sagt, $750$ ml reichen. Wer hat recht? Begründe. \hfill (P10 2015 OS)
\teil Bei einer Gartenbank werden Sitzfläche ($0{,}64$ m²) und Lehne ($0{,}36$ m²) auf beiden Seiten zweimal gestrichen; $1$ l Farbe reicht für $9$ m². Es gibt Dosen mit $250$ ml, $500$ ml und $1$ l. Jan will die 500-ml-Dose, Eva meint, man braucht die 1-l-Dose. Wer hat recht? Begründe. \hfill (P10 2015 OS)
\teil Ein Busunternehmen hat auf $120\,000$ km Mehrkosten von $744$ € für Reifen. In der Zeitung steht: „Jeder Kilometer wird dadurch etwa $6$ ct teurer.“ Die Rechnung dazu: $744$ € $= 744\,000$ ct; $744\,000 : 120\,000 = 6{,}2$ ct. Erkläre den Fehler und korrigiere die Behauptung. \hfill (P10 2014 OS)
\teil Ein Lieferdienst zahlt auf $90\,000$ km $396$ € mehr für Strom. Ein Fahrer behauptet: „Pro Kilometer sind das etwa $4$ ct mehr.“ Seine Rechnung: $396$ € $= 396\,000$ ct; $396\,000 : 90\,000 = 4{,}4$ ct. Erkläre den Fehler und korrigiere die Behauptung. \hfill (P10 2014 OS)
\end{teile}
\end{aufgabe}
